\documentclass[../../main.tex]{subfiles}% 注意这里的文件路径不能用 ./main.tex ,否则用latexmk编译子文件会报错
\graphicspath{{\subfix{./image/}}{\subfix{../../image/}}} % 指定图片引用路径,后续可以直接使用图片文件名
% 如果图片存放在上级目录,则使用 ../../image/ 作为备用路径

% 例如:
% \begin{figure}[H]
% \centering
% \includegraphics[width=0.8\textwidth]{图.png}
% \caption{}
% \label{figure:图}
% \end{figure}
% 注意:上述\label{}一定要放在\caption{}之后,否则引用图片序号会只会显示??.

\begin{document}

\section{曲面上的曲线}

\begin{definition}
设曲面$S$上给定一阶标架场$\{\boldsymbol{r};\boldsymbol{\alpha}_1,\boldsymbol{\alpha}_2,\boldsymbol{\alpha}_3\}$，其中$\boldsymbol{\alpha}_3=\boldsymbol{n}$，其相对分量满足$\omega^3=0,\ \omega^i_j=-\omega^j_i$。对$S$上以弧长$s$为参数的曲线$C:u^\alpha=u^\alpha(s)$，取单位正交标架场
\begin{align*}
\boldsymbol{e}_1=\cos\theta\boldsymbol{\alpha}_1+\sin\theta\boldsymbol{\alpha}_2,\quad \boldsymbol{e}_2=-\sin\theta\boldsymbol{\alpha}_1+\cos\theta\boldsymbol{\alpha}_2,\quad \boldsymbol{e}_3=\boldsymbol{\alpha}_3=\boldsymbol{n},
\end{align*}
称$\{\boldsymbol{r};\boldsymbol{e}_1,\boldsymbol{e}_2,\boldsymbol{e}_3\}$为该一阶标架场沿曲线$C$的限制，其中$\theta$为$\boldsymbol{e}_1$与$\boldsymbol{\alpha}_1$的方向角。

若一阶标架场$\{\boldsymbol{r};\boldsymbol{\alpha}_1,\boldsymbol{\alpha}_2,\boldsymbol{\alpha}_3\}$中$\boldsymbol{\alpha}_1,\boldsymbol{\alpha}_2$沿曲面的主方向，则称其为曲面$S$的\textbf{二阶标架场}。
\end{definition}

\begin{proposition}
设$\theta$为曲面$S$上曲线$C$的切向与$\boldsymbol{\alpha}_1$的方向角，则曲线$C$的测地曲率、法曲率、测地挠率分别为
\begin{align}
&\kappa_g=\frac{\mathrm{d}\theta}{\mathrm{d}s}+\frac{\omega_1^2}{\mathrm{d}s},\label{eq:geodesic-curvature}\\
&\kappa_n=\frac{\omega^1\omega_1^3+\omega^2\omega_2^3}{\mathrm{d}s^2}=a\cos^2\theta+2b\sin\theta\cos\theta+c\sin^2\theta,\label{eq:normal-curvature}\\
&\tau_g=\frac{\omega^1\omega_2^3-\omega^2\omega_1^3}{\mathrm{d}s^2}=b\cos^2\theta+(c-a)\cos\theta\sin\theta-b\sin^2\theta,\label{eq:geodesic-torsion}
\end{align}
其中$a,b,c$满足$(\omega_1^3,\omega_2^3)=(\omega^1,\omega^2)\begin{pmatrix}a&b\\b&c\end{pmatrix}$。
\end{proposition}

\begin{proof}
由单位切向量$\boldsymbol{e}_1=\frac{\mathrm{d}\boldsymbol{r}}{\mathrm{d}s}=\frac{\omega^1}{\mathrm{d}s}\boldsymbol{\alpha}_1+\frac{\omega^2}{\mathrm{d}s}\boldsymbol{\alpha}_2=\cos\theta\boldsymbol{\alpha}_1+\sin\theta\boldsymbol{\alpha}_2$，可得$\frac{\omega^1}{\mathrm{d}s}=\cos\theta,\ \frac{\omega^2}{\mathrm{d}s}=\sin\theta$。
计算曲率向量：
\begin{align*}
\frac{\mathrm{d}\boldsymbol{e}_1}{\mathrm{d}s}=\left(\frac{\mathrm{d}\theta}{\mathrm{d}s}+\frac{\omega_1^2}{\mathrm{d}s}\right)\boldsymbol{e}_2+\frac{\omega^1\omega_1^3+\omega^2\omega_2^3}{\mathrm{d}s^2}\boldsymbol{e}_3.
\end{align*}
测地曲率为曲率向量的切向分量，法曲率为曲率向量的法向分量，即得式\eqref{eq:geodesic-curvature}与\eqref{eq:normal-curvature}。
再计算
\begin{align*}
\frac{\mathrm{d}\boldsymbol{e}_2}{\mathrm{d}s}=-\left(\frac{\mathrm{d}\theta}{\mathrm{d}s}+\frac{\omega_1^2}{\mathrm{d}s}\right)\boldsymbol{e}_1+\frac{\omega^1\omega_2^3-\omega^2\omega_1^3}{\mathrm{d}s^2}\boldsymbol{e}_3,
\end{align*}
测地挠率为其法向分量，即得式\eqref{eq:geodesic-torsion}。

\end{proof}

\begin{theorem}
曲面一点处法曲率$\kappa_n$的极值为该点\textbf{主曲率}$\kappa_1,\kappa_2$，满足
\begin{align}
\kappa_1=\frac{a+c}{2}+\sqrt{\left(\frac{a-c}{2}\right)^2+b^2},\quad \kappa_2=\frac{a+c}{2}-\sqrt{\left(\frac{a-c}{2}\right)^2+b^2},\label{eq:principal-curvature}\\
2H=\kappa_1+\kappa_2=a+c,\quad K=\kappa_1\kappa_2=ac-b^2,\label{eq:mean-gauss-curvature}
\end{align}
其中$H$为平均曲率，$K$为Gauss曲率；且法曲率满足\textbf{Euler公式}
\begin{align}
\kappa_n=\kappa_1\cos^2\theta+\kappa_2\sin^2\theta,\label{eq:euler-formula}
\end{align}
$\theta$为切方向与主方向的夹角。
\end{theorem}

\begin{proof}
将法曲率三角展开：
\begin{align*}
\kappa_n=\frac{a+c}{2}+\frac{a-c}{2}\cos2\theta+b\sin2\theta.
\end{align*}
取$\theta_0$满足
\begin{align*}
\cos2\theta_0=\frac{a-c}{\sqrt{(a-c)^2+4b^2}},\quad \sin2\theta_0=\frac{2b}{\sqrt{(a-c)^2+4b^2}},
\end{align*}
则
\begin{align*}
\kappa_n=\frac{a+c}{2}+\sqrt{\left(\frac{a-c}{2}\right)^2+b^2}\cos2(\theta-\theta_0).
\end{align*}
故$\kappa_n$在$\theta=\theta_0,\theta_0+\pi$取最大值$\kappa_1$，在$\theta=\theta_0+\frac{\pi}{2},\theta_0+\frac{3\pi}{2}$取最小值$\kappa_2$，即得主曲率。
直接计算可得平均曲率与Gauss曲率满足式\eqref{eq:mean-gauss-curvature}。
代入化简即得Euler公式\eqref{eq:euler-formula}。

\end{proof}

\begin{proposition}
曲面$S$上曲率线的微分方程为
\begin{align}
\omega^1\omega_2^3-\omega^2\omega_1^3=0,\label{eq:curvature-line-eq}
\end{align}
且法曲率与测地挠率满足$\frac{\mathrm{d}\kappa_n(\theta)}{\mathrm{d}\theta}=2\tau_g(\theta)$。
\end{proposition}

\begin{proof}
主方向处测地挠率$\tau_g=0$，即$\omega^1\omega_2^3-\omega^2\omega_1^3=0$，故该式为曲率线微分方程。
由$\tau_g=\frac{1}{2}(\kappa_2-\kappa_1)\sin2(\theta-\theta_0)$，对$\kappa_n$求导可得$\frac{\mathrm{d}\kappa_n(\theta)}{\mathrm{d}\theta}=2\tau_g(\theta)$。

\end{proof}

\begin{proposition}
在曲面$S$的二阶标架场下，有$\omega_1^3=\kappa_1\omega^1,\ \omega_2^3=\kappa_2\omega^2$，第二基本形式简化为
\begin{align}
\mathrm{II}=\kappa_1(\omega^1)^2+\kappa_2(\omega^2)^2.\label{eq:ii-second-order-frame}
\end{align}
\end{proposition}

\begin{proof}
二阶标架场中$\theta_0=0$，故$b=0$，此时$a=\kappa_1,\ c=\kappa_2$，代入$(\omega_1^3,\omega_2^3)=(\omega^1,\omega^2)\begin{pmatrix}a&0\\0&c\end{pmatrix}$，得$\omega_1^3=\kappa_1\omega^1,\ \omega_2^3=\kappa_2\omega^2$，代入第二基本形式表达式即得式\eqref{eq:ii-second-order-frame}。

\end{proof}

\begin{example}
在正交参数系$(u,v)$下，曲面第一基本形式为$\mathrm{I}=E(\mathrm{d}u)^2+G(\mathrm{d}v)^2$，推导测地曲率的Liouville公式。
\end{example}

\begin{solution}

取一阶标架场相对分量$\omega^1=\sqrt{E}\mathrm{d}u,\ \omega^2=\sqrt{G}\mathrm{d}v$，联络形式$\omega_1^2=-\frac{(\sqrt{E})_v}{\sqrt{G}}\mathrm{d}u+\frac{(\sqrt{G})_u}{\sqrt{E}}\mathrm{d}v$。
由$\cos\theta=\sqrt{E}\frac{\mathrm{d}u}{\mathrm{d}s},\ \sin\theta=\sqrt{G}\frac{\mathrm{d}v}{\mathrm{d}s}$，代入式\eqref{eq:geodesic-curvature}得
\begin{align*}
\kappa_g&=\frac{\mathrm{d}\theta}{\mathrm{d}s}-\frac{(\sqrt{E})_v}{\sqrt{G}}\frac{\mathrm{d}u}{\mathrm{d}s}+\frac{(\sqrt{G})_u}{\sqrt{E}}\frac{\mathrm{d}v}{\mathrm{d}s}\\
&=\frac{\mathrm{d}\theta}{\mathrm{d}s}-\frac{1}{2\sqrt{G}}\frac{\partial\ln E}{\partial v}\cos\theta+\frac{1}{2\sqrt{E}}\frac{\partial\ln G}{\partial u}\sin\theta.
\end{align*}

\end{solution}

\begin{example}
设曲面$S$的第一基本形式为$\mathrm{I}=(\mathrm{d}u)^2+(u^2+a^2)(\mathrm{d}v)^2$，求曲线$C:u+v=0$的测地曲率。
\end{example}

\begin{solution}

作参数变换$\xi=\frac{1}{2}(u-v),\ \eta=\frac{1}{2}(u+v)$，则$u=\xi+\eta,\ v=-\xi+\eta$，曲线$C$对应$\eta=0$。
代入第一基本形式并配平方，取一阶标架场相对分量
\begin{align*}
\omega^1&=\sqrt{1+a^2+(\xi+\eta)^2}\mathrm{d}\xi+\frac{1-a^2-(\xi+\eta)^2}{\sqrt{1+a^2+(\xi+\eta)^2}}\mathrm{d}\eta,\\
\omega^2&=\frac{2\sqrt{a^2+(\xi+\eta)^2}}{\sqrt{1+a^2+(\xi+\eta)^2}}\mathrm{d}\eta.
\end{align*}
此时$\eta=$常数的曲线满足$\omega^2=0$，$\boldsymbol{\alpha}_1$为$\xi-$曲线切向，限制在$C$上时$\boldsymbol{\alpha}_1$为$C$的切向，故$\theta=0$，$\kappa_g=\frac{\omega_1^2}{\mathrm{d}s}\big|_{\eta=0}$。
设$\omega_1^2=p\mathrm{d}\xi+q\mathrm{d}\eta$，则$\kappa_g=p\cdot\frac{\mathrm{d}\xi}{\mathrm{d}s}\big|_{\eta=0}$。
由结构方程$\mathrm{d}\omega^1=\omega^2\wedge\omega_1^2$，计算外微分：
\begin{align*}
\mathrm{d}\omega^1=\left(\frac{-3(\xi+\eta)}{\sqrt{1+a^2+(\xi+\eta)^2}}-\frac{(1-a^2-(\xi+\eta)^2)(\xi+\eta)}{\left(\sqrt{1+a^2+(\xi+\eta)^2}\right)^3}\right)\mathrm{d}\xi\wedge\mathrm{d}\eta,
\end{align*}
代入$\eta=0$得
\begin{align*}
p\big|_{\eta=0}=\frac{-\xi(2+a^2+\xi^2)}{\sqrt{a^2+\xi^2}(1+a^2+\xi^2)}.
\end{align*}
曲线$C$的弧长元$\mathrm{d}s^2=(1+a^2+\xi^2)\mathrm{d}\xi^2$，故$\frac{\mathrm{d}\xi}{\mathrm{d}s}\big|_{\eta=0}=\frac{1}{\sqrt{1+a^2+\xi^2}}$。
因此
\begin{align*}
\kappa_g=\frac{-\xi(2+a^2+\xi^2)}{\sqrt{a^2+\xi^2}\sqrt{(1+a^2+\xi^2)^3}}=-\frac{u(2+a^2+u^2)}{\sqrt{a^2+u^2}\sqrt{(1+a^2+u^2)^3}},
\end{align*}
其中利用$u+v=0$即$\xi=u$化简。

\end{solution}

\begin{theorem}
设$S$为光滑有向曲面，$C$是$S$上连续可微的简单闭曲线，其围成$S$内的单连通区域$D$，$C$取区域$D$诱导的定向。记$K$为曲面$S$的\textbf{Gauss曲率}，$\kappa_g$为曲线$C$的\textbf{测地曲率}，$s$为$C$的弧长参数，则成立\textbf{Gauss-Bonnet公式}:
\begin{align}
\oint_C \kappa_g \mathrm{d}s + \iint_D K \omega^1 \land \omega^2 = 2\pi. \label{eq:gauss_bonnet_single}
\end{align}
\end{theorem}

\begin{proof}
取区域$D$上的单位正交切向量场$\{\boldsymbol{e}_1,\boldsymbol{e}_2\}$，构成$D$上的一阶标架场，$\omega^1,\omega^2$为标架的相对分量，$\omega_1^2=-\omega_2^1$为对应的联络形式。
设$\theta$为$C$的单位切向量$\boldsymbol{\alpha}$与标架向量$\boldsymbol{e}_1$的夹角，则测地曲率满足
\begin{align*}
\kappa_g = \frac{\mathrm{d}\theta}{\mathrm{d}s} + \frac{\omega_1^2}{\mathrm{d}s},
\end{align*}
整理得
\begin{align}
\mathrm{d}\theta = \kappa_g \mathrm{d}s - \omega_1^2. \label{eq:geo_curv_diff}
\end{align}

因$C$紧致，可取连续可微的方向角分支$\theta(s)$，沿闭曲线$C$对\eqref{eq:geo_curv_diff}积分：
\begin{align*}
\oint_C \mathrm{d}\theta &= \theta(L)-\theta(0) = \oint_C \kappa_g \mathrm{d}s - \oint_C \omega_1^2 \\
&= \oint_C \kappa_g \mathrm{d}s - \iint_D \mathrm{d}\omega_1^2 \\
&= \oint_C \kappa_g \mathrm{d}s + \iint_D K \omega^1 \land \omega^2,
\end{align*}
其中第三步使用Stokes公式，第四步代入Gauss曲率表达式$K=-\frac{\mathrm{d}\omega_1^2}{\omega^1\land\omega^2}$。

在$D$内取定点$p$，以$p$为中心作半径$\varepsilon$的测地圆周$C_\varepsilon$，围成区域$D_\varepsilon\subset D$，取诱导定向，则$\partial(D\setminus D_\varepsilon)=C-C_\varepsilon$。对\eqref{eq:geo_curv_diff}在$\partial(D\setminus D_\varepsilon)$上积分：
\begin{align*}
\int_{C-C_\varepsilon} \mathrm{d}\theta &= \int_{C-C_\varepsilon} \kappa_g \mathrm{d}s - \iint_{D\setminus D_\varepsilon} \mathrm{d}\omega_1^2 \\
&= \int_{C-C_\varepsilon} \kappa_g \mathrm{d}s + \iint_{D\setminus D_\varepsilon} K \omega^1 \land \omega^2.
\end{align*}

在$D\setminus D_\varepsilon$上取特殊一阶标架场$\{\boldsymbol{\alpha}_1,\boldsymbol{\alpha}_2\}$，使$\boldsymbol{\alpha}_1$在$C,C_\varepsilon$上分别为其切向量，此时方向角$\theta\equiv0$，代入得
\begin{align}
\int_{C-C_\varepsilon} \kappa_g \mathrm{d}s + \iint_{D\setminus D_\varepsilon} K \omega^1 \land \omega^2 = 0. \label{eq:zero_int}
\end{align}

因此有
\begin{align}
\oint_C \mathrm{d}\theta = \oint_{C_\varepsilon} \mathrm{d}\theta. \label{eq:angle_eq}
\end{align}

测地圆周$C_\varepsilon$绕行一周时，切向量方向角总变差为$2\pi$，即$\oint_{C_\varepsilon}\mathrm{d}\theta=2\pi$，故$\theta(L)-\theta(0)=2\pi$，代入即得\eqref{eq:gauss_bonnet_single}.

\end{proof}

\begin{example}
设曲面$S$的第一基本形式是
\begin{align*}
\mathrm{I}=(\mathrm{d}u)^2+2\cos\psi\mathrm{d}u\mathrm{d}v+(\mathrm{d}v)^2,
\end{align*}
其中$\psi$是$u$-曲线和$v$-曲线的夹角。试求:
\begin{enumerate}[(1)]
\item 曲面$S$的Gauss曲率$K$;
\item $u$-曲线和$v$-曲线的测地曲率;
\item $u$-曲线和$v$-曲线的二等分角轨线的测地曲率.
\end{enumerate}
\end{example}

\begin{example}
设曲面$S$的第一基本形式是
\begin{align*}
\mathrm{I}=\frac{a^2}{v^2}\left((\mathrm{d}u)^2+(\mathrm{d}v)^2\right),
\end{align*}
求曲线$u=kv+b$($k,b$是常数)的测地曲率.
\end{example}

\begin{example}
设曲线$C$是欧氏空间$E^3$中的一条曲线，其曲率处处不是零。以曲线$C$上每一点为中心在曲线的法平面内作半径为$a$的圆周，这些圆周生成一个曲面，成为围绕曲线$C$的管状曲面，记为$S_a$。试在管状曲面$S_a$上建立一个一阶标架场，写出曲面$S_a$的第一基本形式和第二基本形式，并且求管状曲面$S_a$的主方向和主曲率.
\end{example}





\end{document}